\DocumentMetadata{
  pdfversion=2.0,
  pdfstandard=ua-2,
  lang=en-US,
  tagging=on,
  tagging-setup={math/setup=mathml-SE,tabsorder=structure}
}
\documentclass[11pt,letterpaper]{article}
\usepackage[margin=0.68in,includefoot]{geometry}
\usepackage{newcomputermodern}
\usepackage{amsmath,amscd,mathtools}
\providecommand{\square}{\mdlgwhtsquare}
\usepackage[hidelinks]{hyperref}
\hypersetup{
  pdftitle={Combinatorics PhD Exam, May 2018},
  pdfauthor={University of Florida Department of Mathematics},
  pdfsubject={Graduate mathematics examination},
  pdfdisplaydoctitle=true,
  bookmarksopen=true
}
\setcounter{secnumdepth}{0}
\setlength{\parindent}{0pt}
\setlength{\parskip}{0.28em}
\setlength{\emergencystretch}{3em}
\setlength{\leftmargini}{2.15em}
\setlength{\leftmarginii}{2.65em}
\setlength{\leftmarginiii}{3em}
\renewcommand{\labelenumi}{\textbf{\theenumi.}}
\pagestyle{empty}
\raggedbottom
\allowdisplaybreaks
\newenvironment{examproblems}[1]{%
  \begin{enumerate}%
  \setcounter{enumi}{#1}%
  \setlength{\topsep}{0.35em}%
  \setlength{\partopsep}{0pt}%
  \setlength{\itemsep}{0.48em}%
  \setlength{\parsep}{0.18em}%
}{\end{enumerate}}
\newenvironment{examsubparts}{%
  \begin{enumerate}%
  \setlength{\topsep}{0.18em}%
  \setlength{\partopsep}{0pt}%
  \setlength{\itemsep}{0.16em}%
  \setlength{\parsep}{0.1em}%
}{\end{enumerate}}
\newenvironment{examinstructions}{%
  \par\begingroup\itshape
}{%
  \par\endgroup\vspace{0.18em}
}
\makeatletter
\renewcommand\section{\@startsection{section}{1}{0pt}%
  {-1.25ex plus -.25ex minus -.1ex}%
  {0.55ex plus .1ex}%
  {\normalfont\large\bfseries}}
\renewcommand{\@maketitle}{%
  \newpage\null\vskip -1.2em%
  {\footnotesize\bfseries
    \href{https://www.ufl.edu/}{University of Florida}, \href{https://math.ufl.edu/}{Department of Mathematics}\par}%
  \vskip 0.18em%
  \tagstructbegin{tag=Title}%
    {\LARGE\bfseries\href{https://gma.math.ufl.edu/exams/combinatorics-phd/}{Combinatorics PhD Exam}, May 2018\par}%
  \tagstructend%
  \vskip 0.35em%
  \tagmcbegin{artifact}%
    \hrule height 1pt%
  \tagmcend%
  \vskip 0.55em%
}
\makeatother
\title{Combinatorics PhD Exam, May 2018}
\author{}
\date{}
\begin{document}
\maketitle
\thispagestyle{empty}
\begin{examproblems}{0}
\item
Consider the set \(T_n\) of (not rooted) trees with \(n \ge 3\) labeled leaves for which each interior (non-leaf) vertex has degree 3.
\begin{examsubparts}
\item[\textnormal{(a)}]
Prove that each tree in \(T_n\) has exactly \(2n-3\) edges.
\item[\textnormal{(b)}]
Prove that the number of trees in \(T_n\) is \((2n-5)!!=1\cdot 3\cdot 5\cdot 7\cdots(2n-7)\cdot(2n-5)\).
\end{examsubparts}
\item
The famous Dilworth Theorem for a finite poset~\(P\) states that the minimum number of chains in any partition of~\(P\) into chains is equal to the maximum number of elements in an antichain of~\(P\).
\begin{examsubparts}
\item[\textnormal{(a)}]
Prove that if~\(P\) has size at least \(rs+1\), then there is either a chain of size~\(r\) or an antichain of size~\(s\).
\item[\textnormal{(b)}]
There are many proofs of the following Erdős–Szekeres theorem: For any sequence \(a_1,a_2,\ldots,a_{n^2+1}\) of integers, there is a subsequence of length \(n+1\) that is monotone. Prove the Erdős–Szekeres theorem using Dilworth’s theorem, i.e., part (a).
\end{examsubparts}
\item
Prove that 
\[
\sum_{i=0}^{n}(-1)^i\binom{n}{i}i^k=
\begin{cases}
0 & \text{if }0\le k<n,\\
(-1)^n n! & \text{if }k=n.
\end{cases}
\]
 What is the sum if \(k>n\)?

\par
Hint: Count surjective mappings from \([k]\) to \([n]\).
\item
Recall: The girth of a graph is the length of the shortest cycle (and~\(\infty\) if the graph has no cycles).
\begin{examsubparts}
\item[\textnormal{(a)}]
Prove that a planar graph of girth (minimum cycle length) at least 6 has a vertex of degree at most 2.
\item[\textnormal{(b)}]
The Four Color Theorem states that any planar graph is 4-colorable. Grötzsch proved that any triangle-free planar graph is 3-colorable. Without using his result, prove that a planar graph of girth at least 6 is 3-colorable.
\end{examsubparts}
\item
How many strings can be formed using the alphabet \(\{A,B,C,D,E\}\) if
\begin{examsubparts}
\item[\textnormal{(a)}]
the letter~\(A\) occurs an odd number of times
\item[\textnormal{(b)}]
the letters~\(A\) and~\(B\) are both used an odd number of times.
\end{examsubparts}
\item
Let \(f(n)\) be the number of ways to tile a \(1\times n\) path with \(1\times 2\) tiles that are red, blue, or green, and \(1\times 1\) tiles that are yellow, orange, black, or white.
\begin{examsubparts}
\item[\textnormal{(a)}]
Find an explicit formula for \(f(n)\).
\item[\textnormal{(b)}]
On average, how many \(1\times 1\) tiles will a \(1\times n\) path contain?
\end{examsubparts}
\item
Choose a derangement~\(p\) of length~\(n\) uniformly at random. Recall that a derangement is a permutation with no 1-cycles. On average, how many 2-cycles does~\(p\) contain?
\item
Let~\(C\) be a binary code of length~\(n\) and minimum distance \(d\ge 2e+1\). Prove the following two bounds (the Hamming bound and Singleton bound). 
\[
|C|\le \frac{2^n}{\displaystyle\sum_{i=0}^{e}\binom{n}{i}}
\]
 
\[
|C|\le 2^{n-d+1}
\]
\end{examproblems}
\end{document}
